\chapter{Conclusion and Future Work}\label{ch:conclusion}

\section{What we built and what we found}
TrustLens is an evidence-driven, FRIES-aligned model auditing framework. It imports a Hugging Face text-classification model, runs five probes locally, stores hashed evidence, interprets it into Occurrence, Severity and Detection ratings, and produces a versioned report whose primary output is a five-dimensional audit profile. Abstention is a first-class result, and the aggregate FRIES score is a secondary summary.

We tested the system in ten validation studies; Table~\ref{tab:rq} gives the answer to each research question. The measured evidence is reproducible and rises monotonically with injected severity ($\rho=1.0$ in every sweep). Calibrating the O/S/D mapping widened the score range over the injected defects on held-out sweeps from 0.68--1.18 to 4.60--6.65 points (a scale relative to those defects). Separating disclosure gaps from risks removed the false flags on clean controls. The final equal-opportunity rule detected every injected fairness defect at flip rates of 0.4 and above on fresh data, although it also flagged the real reference model. Artifact integrity detected a one-byte change in a demonstration, and TrustLens produced a profile for each of seven real Hub models. Independent reviewers agree with TrustLens on risk for 6 of 12 packets and on sufficiency for 9 of 12.

\section{Objectives and results}
\begin{table}[!htbp]
\centering\footnotesize
\caption{Objectives of Chapter~\ref{ch:intro} and where they are met.}
\begin{tabularx}{\textwidth}{@{}p{4.6cm}L@{}}
\toprule
\textbf{Objective / question} & \textbf{Outcome}\\
\midrule
O1: working audit pipeline & Import, local inference, five probes, O/S/D, scoring, review and report (Chapters~\ref{ch:arch}--\ref{ch:impl}); 913 backend tests (Chapter~\ref{ch:testing}).\\
O2: honest handling of evidence & Abstention, disclosure gaps separated from risks, methodology versions stored (Chapter~\ref{ch:fries}).\\
O3: tested on controlled, reviewed and real models & Chapter~\ref{ch:results}.\\
RQ1--RQ6 & Answered in Table~\ref{tab:rq}.\\
\bottomrule
\end{tabularx}
\end{table}

\medskip\noindent Beyond the results (scope in Chapter~\ref{ch:limits}), we hope the method carries over: separate the layers, declare before running, freeze the folders, report the failures.

\section{Future work}
\begin{enumerate}[nosep]
\item \textbf{Validate the decision layer on new data}, including the v10 safety rule.
\item \textbf{A larger reviewer study} with distinct identifiers, domain experts and an ethics review.
\item \textbf{Repeat the LLM comparison with one pinned provider}, ideally local.
\item \textbf{More than one risk per dimension}, as the FRIES procedure intends.
\item \textbf{Broader probes:} gradient-guided perturbations, SHAP or LIME, intersectional groups, other tasks, signed manifests.
\end{enumerate}
